% =============================================================================
% Appendix: knowledge-area annotation --- method, taxonomy evolution and final
% distribution.
%
% Usage: \input this file after \appendix in the paper.
% Required packages: graphicx, booktabs, amsmath, amssymb, natbib (\citep),
% hyperref or url (\url in footnotes; the ACL template loads hyperref).
% Assumes UTF-8 input and T1 font encoding (both provided by the ACL template).
% Bibliography entries for every \citep below are in
% ablations/knowledge_area/references_knowledge_area.bib.
%
% The annotation method (model, prompt template, decoding, two passes and
% adjudication) was fixed before the first full run and never varied. What
% varied is the taxonomy: versions 1.0 to 1.4, three full runs (1.1, 1.2 and
% 1.4) over the same questions and two runs of 1.3 and 1.4 on the abstentions
% of run 1.2. Every number below comes from the artifacts of those runs under
% output/knowledge-annotation-work/ (full-run, full-run-2,
% uncertain-taxonomy-v1.3, uncertain-taxonomy-v1.4, full-run-taxonomy-v1.4 and
% its post-annotation/ directory), from the taxonomy files in
% src/mmlu_pt/annotation/knowledge_area/ (1.0 is the git snapshot kept in
% ablations/knowledge_area/assets/) and from the versioned lists in config/.
% ablations/knowledge_area/make_figures.py prints every aggregate quoted
% below, checks the exports against the QC reports, the checkpoints and the
% taxonomies, and generates the figures. The only numbers not printed by the
% script are the sizes of the audits that motivated each taxonomy version,
% taken from src/mmlu_pt/annotation/knowledge_area/TAXONOMY_CHANGELOG.md.
% By default figures are read from ./figures; to change this, define the macro
% before the \input, e.g.
%   \newcommand{\KnowledgeAreaFigDir}{ablations/knowledge_area/figures}
% =============================================================================
\providecommand{\KnowledgeAreaFigDir}{figures}

\section{Knowledge-Area Annotation: Method, Taxonomy Evolution and Final Distribution}
\label{app:ka}

Every released question carries a \emph{subject}, chosen from a closed list of
candidates fixed per exam, and a \emph{macro area}, derived from the subject
by a table and never generated. The paper uses them to report accuracy per
area and to stratify the few-shot demonstrations by macro area and academic
level. The labels are assigned by an open-weight language model that reads
the question, its choices, the exam and the edition, chooses one candidate or
abstains with \texttt{UNCERTAIN}, and does so twice with independent seeds; a
third request by the same model settles the disagreements. Unlike the
filtering and deduplication stages, the method itself was not ablated: model,
prompt template, decoding and aggregation were fixed before the first full run
and never changed. What changed was the \emph{taxonomy}. Five versions were
written (1.0 to 1.4), three of them were run over the whole corpus (1.1, 1.2
and 1.4) and two were run on the abstentions of the previous full run (1.3 and
1.4), so the version is the only treatment this appendix can compare
(Table~\ref{tab:ka-overview}). The appendix fixes the notation
(\S\ref{app:ka-setup}), records the method and the reasons behind the choice
of model (\S\ref{app:ka-method}), traces the taxonomy from 1.0 to 1.4
(\S\ref{app:ka-taxonomy}), measures what each version changed in the labels
(\S\ref{app:ka-effect}), describes the final distribution
(\S\ref{app:ka-distribution}), the manual post-annotation
(\S\ref{app:ka-post}) and the cost (\S\ref{app:ka-cost}). Limitations are
discussed in \S\ref{app:ka-limitations}.

\begin{table}[t]
\centering
\footnotesize
\setlength{\tabcolsep}{3.5pt}
\begin{tabular}{@{}lrrrrlrrrr@{}}
\toprule
Ver. & Subj. & Areas & Exams & Rules & Run & $n$ & Agr. & Adj. & Unc. \\
\midrule
1.0 & 70 & 11 & 15 & 8 & -- & -- & -- & -- & -- \\
1.1 & 70 & 11 & 17 & 8 & full & 41,653 & 91.64 & 8.36 & 88 \\
1.2 & 78 & 9 & 17 & 14 & full & 41,653 & 92.14 & 7.86 & 23 \\
1.3 & 78 & 9 & 17 & 15 & subset & 23 & -- & -- & 9 \\
1.4 & 78 & 9 & 17 & 15 & subset & 15 & -- & -- & 1 \\
1.4 & 78 & 9 & 17 & 15 & full & 41,636 & 92.08 & 7.92 & 8 \\
\bottomrule
\end{tabular}
\caption{Taxonomy versions and the runs made with each. \emph{Subj.}, \emph{Areas}, \emph{Exams}, \emph{Rules}: subjects, macro areas, exams with a candidate list and annotation rules in the version; \emph{Run}: a full run over the corpus or a run over the abstentions of the previous full run; $n$: questions; \emph{Agr.}: percentage of questions on which the two independent passes chose the same subject; \emph{Adj.}: percentage sent to adjudication; \emph{Unc.}: questions left \texttt{UNCERTAIN}. Runs 1.1 and 1.2 annotated the 41,653 questions of the filtered corpus; run 1.4 annotated the 41,636 questions left after the dataset revision. Agreement is the self-consistency of one model, not accuracy.}
\label{tab:ka-overview}
\end{table}

\paragraph{Summary.} The method is fixed and the taxonomy is the variable.
(1)~One open-weight model, Qwen3.5-122B-A10B in FP8, labels each question
twice with independent seeds under a JSON-schema grammar that admits only the
exam's candidates and \texttt{UNCERTAIN}; a third request by the same model
adjudicates when the passes differ or either abstains. (2)~The taxonomy grew
from 70 subjects in 11 macro areas and 15 exams (1.0) to 78 subjects in 9
areas and 17 exams with 15 annotation rules (1.4), each step motivated by an
audit of the previous run's abstentions, low-confidence labels and targeted
cases. (3)~Abstentions fell from 88 (1.1) to 23 (1.2) to 8 (1.4) and to 0
after six manual assignments and two removals; the eight abstentions of the
final run are questions no earlier run had abstained on. (4)~Between
consecutive full runs 7.96\% (1.1 to 1.2) and 4.46\% (1.2 to 1.4) of the
labels changed, against a within-run disagreement between the two passes of
7.82\% and 7.92\%; 549 questions took one of the eight subjects added in 1.2.
(5)~All nine macro areas of the final dataset have at least 2,472 questions;
all 78 subjects are used, 14 with fewer than 100 questions. (6)~The final run
took 3\,h\,39\,min on one GPU and processed 311.2 million prompt tokens.

% -----------------------------------------------------------------------------
\subsection{Data, notation and metrics}
\label{app:ka-setup}

\paragraph{Data.} Runs 1.1 and 1.2 annotated $D_F$, the 41,653 questions of
the filtered corpus (\texttt{bench-temp/mmlu-pt-filtered}, revision
\texttt{bef2dcfb}\ldots). Run 1.4 annotated $D_R$, the 41,636 questions of the
revised corpus (\texttt{bench-temp-2/mmlu-pt-revised}, revision
\texttt{0e8b69a1}\ldots), which removes 17 questions with an irrecoverable
defect (the dataset revision); the final dataset removes two more and
has 41,634 questions from 17 exams, 9,789 at high-school and 31,845 at
undergraduate level. A question is identified by the SHA-256 of its exam,
edition, number, statement and choices, so the three full runs join exactly
on the 41,636 questions of $D_R$, and the two subset runs are identified
within the full runs by the same key. A record $x=(q,C,a,e,d,\ell)$ has
statement $q$, ordered choices $C$, answer index $a$, exam $e$, edition $d$
and level $\ell$. The model sees $q$, $C$, $e$ and $d$; it never sees $a$.

\paragraph{Taxonomy and labels.} A taxonomy version $v$ defines a set of
subjects $T_v$ with a one-paragraph definition each, a map
$\mu_v\colon T_v\to A_v$ to macro areas, a candidate list $C_v(e)\subseteq
T_v$ per exam and a list of annotation rules $R_v$ that is included verbatim
in every prompt. A label for a question of exam $e$ is an element of
$C_v(e)\cup\{\texttt{UNCERTAIN}\}$. Two independent passes produce $s_1(x)$
and $s_2(x)$; when $s_1\neq s_2$ or either is \texttt{UNCERTAIN}, an
adjudication request produces $s_a(x)$ and the final label is $s(x)=s_a(x)$
(status \emph{adjudicated}, or \emph{uncertain} when $s_a=\texttt{UNCERTAIN}$);
otherwise $s(x)=s_1(x)$ (status \emph{accepted}) with confidence equal to the
smaller of the two reported confidences. The stored macro area is
$\mu_v(s(x))$. Each pass also returns an alternative subject, a confidence in
$[0,1]$ and a justification of at most 400 characters; the confidence is
model-reported and uncalibrated.

\paragraph{Metrics.} For a run, the \emph{agreement} is the share of
questions with $s_1=s_2\neq\texttt{UNCERTAIN}$, the \emph{adjudication rate}
the share sent to adjudication and the \emph{abstention} the number of
questions with $s=\texttt{UNCERTAIN}$. For two versions $u$ and $v$ run over
the same questions, the \emph{change share} of a group $G$ is
$|\{x\in G: s^{u}(x)\neq s^{v}(x)\}|/|G|$, which is compared with the
within-run \emph{disagreement} $|\{x\in G: s^{v}_1(x)\neq s^{v}_2(x)\}|/|G|$,
the share of labels that one version alone already fails to reproduce between
its two passes. The \emph{uptake} of a version is the number of questions that
received a subject absent from the previous version, or absent from the
previous candidate list of their exam. None of these measures accuracy: no
human-labelled reference exists for these questions, and all three runs use
the same model, so agreement between passes or versions measures consistency,
and a change between versions is not evidence that either label is right.

% -----------------------------------------------------------------------------
\subsection{Annotation method}
\label{app:ka-method}

\paragraph{Motivation.} Three failure modes were anticipated. A model may
label a question by the dominant subject of its exam instead of by its
content; it may choose a neighbouring subject when the right one is missing
from the list; and its labels vary with the prompt and the model, and even
identical inputs can receive different labels under sampling, a variance that
a single sample per item hides \citep{song2025good}. The design answers each: the exam and the edition are given to the
model as context but the prompt states that they are priors, not labels, and
the candidate list is restricted per exam so that an out-of-scope subject
cannot be chosen at all; \texttt{UNCERTAIN} is always available and the rules
forbid substituting an adjacent subject for a missing one; and every question
is labelled twice with independent seeds, in the spirit of self-consistency
checks that sample a model several times \citep{wang2023selfconsistency},
with a third request to settle the cases where the two labels differ.

\paragraph{Model choice.} The annotator is Qwen3.5-122B-A10B-FP8, a
mixture-of-experts model with 122 billion parameters of which 10 billion are
active per token, in the FP8 weights published by its authors (revision
\texttt{a099dee7}\ldots), with its reasoning mode disabled.\footnote{\url{https://huggingface.co/Qwen/Qwen3.5-122B-A10B-FP8}, revision \texttt{a099dee70ccfcd8d5dda56aaa0b60cb8ecadabc9}, Apache 2.0 license. The model card is the source of the parameter counts, the language coverage and the sampling guidance quoted here.} Four considerations led to it; none was tested
against an alternative model on this task. (i)~Large instruction-tuned
models are an established substitute for crowd and expert annotators on
closed-label text classification, including zero-shot classification from a
list of labels with definitions
\citep{ding2023gpt3,he2024annollm,tan2024survey},
with the caveat that their labels must be validated per task
\citep{gligoric2025unconfident}; this appendix therefore reports consistency
rather than accuracy and the stage ends with a manual review
(\S\ref{app:ka-post}). (ii)~The annotator must not be one of the systems the
benchmark evaluates: language models used as judges favour their own outputs,
and the bias grows with their ability to recognize them
\citep{panickssery2024llm,liu2024narcissistic}. All 19 open-weight checkpoints
in the evaluation matrix have at most 35 billion parameters and the
122B-A10B model is not among them, although six of them belong to the same
Qwen3.5 family (\S\ref{app:ka-limitations}). (iii)~Open weights with a pinned
revision make the annotation reproducible and auditable, which an API model
that can change silently does not offer, and open models have been shown to
be viable annotators \citep{alizadeh2025opensource}. The Qwen family
documents broad multilingual coverage, 119 languages including Portuguese for
Qwen3 and 201 languages and dialects for Qwen3.5 according to their technical
report and model card, and is the base of the Portuguese Tucano~2 models
evaluated in the paper. (iv)~The model fits the
available hardware: the FP8 weights occupy 115.37~GiB of the 191.5~GB of one
NVIDIA B200, leaving a KV cache of 1,413,120 tokens, enough for the 50
concurrent sequences of 20,480 tokens used by the run; a single GPU annotates
the corpus in under four hours (\S\ref{app:ka-cost}) while the remaining GPUs
of the evaluation budget stay free. The sampling parameters are the ones the
model card recommends for its thinking mode (temperature~1.0, top-$p$~0.95,
top-$k$~20, presence penalty~1.5); they were kept when the reasoning mode was
switched off, so that the two configurations compared in a pilot differed in
one setting only.

\paragraph{Prompt.} The system message states the task (identify the one
academic discipline whose knowledge is primarily required to determine the
correct answer), forbids solving the question or identifying the answer,
forbids classifying by superficial words or setting, states that exam and
edition are contextual priors and not automatic labels, allows
\texttt{UNCERTAIN} when no candidate fits, declares the question and the
choices untrusted data, and asks for a JSON object with the four fields
above. The user message gives the exam and the edition, the annotation rules
$R_v$ as JSON, the numbered candidates $C_v(e)$ each with its definition, and
the question and its choices as a JSON object. The adjudication prompt is the
same message followed by the two proposals, each reduced to its pass number,
subject and justification, and the instruction to assess them independently
(``they may be wrong''), choosing either, another candidate or
\texttt{UNCERTAIN}. Pass 2 never sees pass 1, and the adjudicator never sees
the confidences. The prompt template has one version
(\texttt{subject\_annotation\_v1}); the rules and the definitions are part of
the taxonomy, so a new taxonomy version changes the prompt content, and the
run manifest records the hash of both.

\paragraph{Decoding and validation.} Inference runs offline in vLLM~0.25.0
\citep{kwon2023efficient} on one GPU, with batches of 128 questions, at most 50 concurrent sequences,
CUDA graphs enabled and prefix caching disabled. Each request is constrained
by a JSON schema whose \texttt{subject} and \texttt{alternative\_subject}
fields are enumerations over $C_v(e)\cup\{\texttt{UNCERTAIN}\}$, enforced by
the XGrammar backend \citep{dong2024xgrammar}, with a budget of 8,192 output
tokens. The parsed object is then validated locally: the subject must be a
candidate or \texttt{UNCERTAIN}, the confidence a number in $[0,1]$, the
justification between 1 and 400 characters, and duplicate keys, NaN values
and reasoning tags are rejected. A request that fails validation is retried
with a new seed, at most five times; the seed of every request is derived by
hashing the base seed (42), the question identifier, the stage, the pass and
the attempt, and is stored with the request. In the final run, 86,571
requests succeeded and six were
retried after the model returned an object that violated the schema
(\texttt{invalid\_final\_schema}); all retries succeeded and every successful
request ended at a stop token rather than at the token budget. Every input,
candidate list, attempt, parsed result and token count is written to a SQLite
checkpoint, and a run resumes only under the same source revision, model
revision, taxonomy, prompt and code hashes.

\paragraph{Aggregation.} The two passes agree on an admissible subject for
92.08\% of the questions in the final run (accepted), and 3,299 questions
(7.92\%) went to adjudication. The adjudicator sided with pass~1 in 1,574
cases, with pass~2 in 1,640, chose a third subject in 77 and abstained in 8;
the earlier full runs show the same pattern (1.1: 1,778, 1,543, 73 and 88;
1.2: 1,561, 1,618, 70 and 23). The near symmetry between the two passes is
expected, since they are exchangeable samples from the same model.

% -----------------------------------------------------------------------------
\subsection{Taxonomy evolution}
\label{app:ka-taxonomy}

\paragraph{Motivation.} The labels must support MMLU-style reporting by
subject while matching what Brazilian exams assess: a school curriculum in
the vestibular, military and olympiad exams, and the professions and degree
programmes of the professional, civil-service and ENADE exams. No official
classification of fields covers both, and the first version was written for
the project. Every later version answers evidence from the previous run:
questions on which the model abstained, labels with low reported confidence,
and cases where the justification named a discipline that the candidate list
did not contain. Table~\ref{tab:ka-versions} lists the versions and
Figure~\ref{fig:ka-candidates} the candidates per exam.

\begin{table*}[t]
\centering
\footnotesize
\setlength{\tabcolsep}{4pt}
\begin{tabular}{@{}lp{0.52\textwidth}p{0.30\textwidth}@{}}
\toprule
Version & Change with respect to the previous version & Evidence \\
\midrule
1.0 & 70 subjects in 11 macro areas, candidate lists for 15 exams, 8 annotation rules. & Written before any run. \\
1.1 & Adds the ENEM (15 candidates) and IME (5) exams; adds Literature to BLUEX (9 to 10) and Public Policy and Public Administration to BNDES (29 to 30). Subjects, areas, definitions and rules unchanged. & Two exams added to the corpus. Pilots on 85 questions, then the first full run. \\
1.2 & Adds 8 subjects (Occupational Therapy; Radiology and Medical Imaging; Aesthetics and Cosmetology; Theology and Religious Studies; Gastronomy and Culinary Arts; Public Safety and Physical Security; Computer Graphics; Human--Computer Interaction); merges 11 macro areas into 9 (Computer Science with Engineering, Architecture and Design; Natural and Earth Sciences with Agricultural and Veterinary Sciences; Social Sciences and Public Policy renamed), moving 24 subjects; widens the candidates of BACEN (18 to 26), BNDES (30 to 36), the residency exams (5 to 11), POSCOMP (9 to 12), BLUEX (10 to 11) and ENADE (70 to 78); refines 29 definitions; rewrites the uncertainty rule and adds 6 rules (compound labels, provided rules, language and metadata, missing context, source integrity, clinical/psychological/legal boundaries). & Audit of run 1.1: its 88 abstentions, 2,901 labels with confidence below 0.90 and targeted cases, 3,011 questions in all, of which 140 received an engineering assessment. Second full run. \\
1.3 & Adds English Language to CNU (50 to 51), Professional Ethics to REVALIDA (5 to 6) and Professional Ethics and Social Security Law to the residency exams (11 to 13); refines 9 definitions; adds the label--justification consistency rule and edits two rules. & Analysis of run 1.2: its 23 abstentions and 48 targeted cases. Run on the 23 abstentions. \\
1.4 & Refines the definition of English Language and edits three rules (language and metadata, source integrity, label--justification consistency). Candidates unchanged. & Run 1.3 on the 23 abstentions (14 labelled, 9 kept); 8 of the 9 removed by the dataset revision. Run on the 15 remaining questions, then the third full run. \\
\bottomrule
\end{tabular}
\caption{Taxonomy versions. \emph{Evidence}: the material that motivated the version and the run made with it. The counts of the audits come from the taxonomy changelog; every other count is computed from the taxonomy files. The assessments that motivated 1.2 and 1.3 are engineering reviews of the model's output against the question text, not human gold labels.}
\label{tab:ka-versions}
\end{table*}

\begin{figure}[t]
\centering
\includegraphics[width=\columnwidth]{\KnowledgeAreaFigDir/taxonomy_versions.pdf}
\caption{Candidate subjects per exam in each taxonomy version. Exams are grouped by academic level and sorted by size in the final dataset; ``--'' marks an exam without a candidate list in that version. ENADE admits every subject of the taxonomy; the vestibular, military and olympiad exams admit the school subjects; the professional exams admit the disciplines of their profession.}
\label{fig:ka-candidates}
\end{figure}

\paragraph{Design.} Four properties hold in every version. The macro area is
a function of the subject and is never offered to the model. Candidate lists
only restrict: an exam-specific list is a subset of the global list, no rule
may depend on the answer, and ENADE, whose booklets cover every degree
programme, admits all subjects rather than a list inferred from the course
named in the edition. Each subject has a definition that names the
neighbouring subjects it excludes, so that the prompt carries the boundaries
(Portuguese Language against Literature, Logic against Mathematics and
Programming, Internal Medicine against General Surgery, the branches of Law
against each other). And \texttt{UNCERTAIN} is a first-class label, so that a
question whose discipline is missing from the list or whose statement is
unreadable is not forced into a neighbour. Law is a macro area of its own with
14 subjects, because the bar exam, the magistracy exam and the civil-service
exams make it the largest field of the corpus; Logic is separate from
Mathematics and from Programming because the informatics olympiad consists
largely of puzzles solved from rules given in the statement.

\paragraph{From 1.0 to 1.1.} Version 1.0 covered the 15 exams of the corpus
at the time; ENEM and IME were added to the corpus afterwards and 1.1 gives
them candidate lists drawn from the school subjects (15 for ENEM, including
Communication Studies for its media items, and 5 for IME). Two candidates
were added to existing lists (Literature to BLUEX, Public Policy and Public
Administration to BNDES); subjects, areas, definitions and rules are
identical, so 1.0 was never run and 1.1 is the baseline.

\paragraph{From 1.1 to 1.2.} The audit of run 1.1 read its 88 abstentions
and 2,901 labels with confidence below 0.90. It found three kinds of problem.
Some disciplines were missing from the taxonomy altogether: radiology items in
the residency and ENADE booklets, occupational therapy, aesthetics,
gastronomy, theology and two computing fields; 1.2 adds the eight subjects.
Some disciplines existed but were absent from the exam's list, so the model
either abstained or chose a neighbour: criminal, civil and labour law in the
BACEN exam, history, geography and sociology in BNDES, psychology, physics and
law in the residency exams, information security in POSCOMP, sociology in
BLUEX; 1.2 widens those lists. And some definitions and rules were
ambiguous, which 1.2 addresses by refining 29 definitions and by adding six
rules, among them that a question whose solution needs only arithmetic on
rules fully supplied in the statement is Mathematics even in a legal or
clinical narrative, that language comprehension is classified by the language
assessed and not by the edition title, and that a missing passage or figure
does not by itself prevent assignment. The merge of 11 macro areas into 9
grouped areas that were too small to report on separately; no subject was
renamed. The longer rules and definitions raised the longest prompt above the
16,384-token context used by run 1.1, so run 1.2 used 20,480 tokens.

\paragraph{From 1.2 to 1.4.} The analysis of run 1.2 read its 23
abstentions and 48 targeted cases. Eight of the 23 were questions whose
statement was unrecoverable (cross-referenced OBI items, a truncated POSCOMP
stem, a BLUEX stem replaced by a formatting instruction), which are a
source-quality problem and were later removed by the dataset revision. The
remaining cases pointed to three candidate gaps (English comprehension in
CNU, professional ethics in the medical exams, social-security law in the
residency exams), to nine definitions and to a new rule requiring the chosen
label to match the discipline named in the justification; the model had, in
some cases, chosen a neighbour while asserting that the right subject was not
in the list. Version 1.3 made those changes and was tested on the 23
abstentions only: 14 received a subject and 9 stayed \texttt{UNCERTAIN}. Of
the nine, eight were the unrecoverable statements; the ninth was an English
item whose passage is missing. Version 1.4 refines the definition of English
Language and three rules so that an item whose stem and choices are entirely
in English is classified as English comprehension even without its passage,
and was tested on the 15 questions left after the revision: 14 received a
subject and one stayed \texttt{UNCERTAIN} (an IME English item whose
identical 2010 copy received English Language in the same run, which shows
sampling noise on an item without its supporting text; both copies were then
removed as well).

\paragraph{Decision.} Version 1.4 was frozen as the default taxonomy before
the third full run. The run manifest records the SHA-256 of the taxonomy
(\texttt{7c43673b}\ldots), of the exam aliases and of the code, and the
figure script checks that the taxonomy stored with each run equals the
packaged file of that version.

% -----------------------------------------------------------------------------
\subsection{Effect of each version on the labels}
\label{app:ka-effect}

\paragraph{Abstentions.} Table~\ref{tab:ka-uncertain} and
Figure~\ref{fig:ka-uncertain} follow the abstentions. Of the 88 questions on
which run 1.1 abstained, 78 received a subject in run 1.2 and 10 did not;
13 questions that 1.1 had labelled became \texttt{UNCERTAIN} in 1.2, for a
total of 23. Run 1.3 on those 23 labelled 14; the dataset revision removed 8
of the 9 it kept; run 1.4 on the 15 remaining labelled 14. The third full
run abstained on 8 questions, none of which any earlier run had abstained on:
two adjudication errors in which the adjudicator asserted that a candidate
was absent from the list, three questions whose discipline is a candidate gap
of the vestibular and residency lists (a Maria da Penha law item in ENEM, a
federative-organization item in COMVEST, a health-regulation item in a
psychiatry residency exam), one boundary case and two questions whose
supporting table or chart is missing. The post-annotation resolved them
(\S\ref{app:ka-post}).

\begin{table}[t]
\centering
\footnotesize
\setlength{\tabcolsep}{4pt}
\begin{tabular}{@{}lrrrr@{}}
\toprule
Stage & Items & Labelled & Uncertain & Removed \\
\midrule
Full run, taxonomy 1.1 & 41,653 & 41,565 & 88 & -- \\
Full run, taxonomy 1.2 & 41,653 & 41,630 & 23 & -- \\
\quad of the 88 abstentions of 1.1 & 88 & 78 & 10 & -- \\
Run on the abstentions, taxonomy 1.3 & 23 & 14 & 9 & -- \\
Dataset revision & 9 & -- & 1 & 8 \\
Run on the remaining, taxonomy 1.4 & 15 & 14 & 1 & -- \\
Full run, taxonomy 1.4 & 41,636 & 41,628 & 8 & -- \\
Post-annotation & 8 & 6 & 0 & 2 \\
\bottomrule
\end{tabular}
\caption{Abstentions at each stage. \emph{Labelled}: questions that received a subject at that stage; \emph{Uncertain}: questions left \texttt{UNCERTAIN}; \emph{Removed}: questions removed from the dataset because their statement is unrecoverable. The 23 abstentions of run 1.2 are the 10 kept from run 1.1 plus 13 new ones; the 8 abstentions of run 1.4 are all new.}
\label{tab:ka-uncertain}
\end{table}

\begin{figure}[t]
\centering
\includegraphics[width=\columnwidth]{\KnowledgeAreaFigDir/uncertain_flow.pdf}
\caption{What became of the abstentions after the first full run. Each bar is the set of questions entering the stage; the first bar is the 88 abstentions of run 1.1 as labelled by run 1.2, and the fifth is the 8 new abstentions of run 1.4.}
\label{fig:ka-uncertain}
\end{figure}

\paragraph{Label changes.} On the 41,636 questions common to the three full
runs, 3,316 labels (7.96\%) changed from 1.1 to 1.2 and 1,855 (4.46\%) from
1.2 to 1.4. The within-run disagreement between the two passes is 7.82\% in
run 1.2 and 7.92\% in run 1.4 on the same questions, so the change from 1.2
to 1.4 is below the level at which one version fails to reproduce its own
labels, and the change from 1.1 to 1.2 is at that level overall but well
above it in three exams whose candidate lists 1.2 widened
(Figure~\ref{fig:ka-changes}, with the disagreement of run 1.4 as reference):
21.56\% of the BACEN labels changed (disagreement 6.45\%), 15.69\% of the
residency labels (11.05\%) and 9.91\% of POSCOMP (5.91\%); ENADE (12.84\%
against 12.69\%) and BNDES (8.97\% against 8.33\%) are at their floor and CNU
(11.20\% against 15.48\%) below it, as are AFA (0.09\%), IME (1.05\%), CFCES
(2.12\%) and FUVEST (2.59\%), whose lists did not change. From 1.2 to 1.4 no
exam exceeds its floor; the largest change share is again CNU (12.83\%), the
one exam whose list changed between those versions (English Language). 907 questions of run 1.2
carry a subject that was not a candidate of their exam in 1.1 (ENADE 406,
BACEN 198, residency 141, POSCOMP 73, BNDES 65, BLUEX 24), and 549 carry one
of the eight new subjects (Radiology and Medical Imaging 167, Gastronomy and
Culinary Arts 99, Computer Graphics 65, Aesthetics and Cosmetology 58,
Occupational Therapy 56, Theology and Religious Studies 54, Public Safety and
Physical Security 25, Human--Computer Interaction 25). The macro area is more
stable than the subject: 92.04\% of the questions keep their subject from
1.1 to 1.2 and 96.38\% keep their macro area once the 1.1 subjects are
projected onto the nine areas of 1.2; from 1.2 to 1.4 the figures are
95.54\% and 97.76\%. Table~\ref{tab:ka-transitions} lists the most frequent
transitions. Those from 1.1 to 1.2 follow the changes of the version: Logic to
Mathematics and Finance, Physics and Statistics to Mathematics follow the
provided-rules rule; Internal Medicine and Physics to Radiology, and Nutrition
to Gastronomy, follow the new subjects. Those from 1.2 to 1.4, whose
definitions differ in one subject, are the boundaries that the two passes
also disagree on within a run (General Surgery and Internal Medicine,
Communication Studies and Portuguese Language, Mathematics and Logic,
Constitutional and Administrative Law), with one exception, Professional Ethics
to Civil Law and Civil Procedure, two subjects whose definitions 1.3 refined.

\begin{figure}[t]
\centering
\includegraphics[width=\columnwidth]{\KnowledgeAreaFigDir/label_changes_by_version.pdf}
\caption{Share of each exam's labels that changed between consecutive full runs, next to the share on which the two passes of run 1.4 disagree. A change share above the disagreement marks an exam whose labels the version moved beyond the sampling noise of the method; the exams clearly above it from 1.1 to 1.2 are among those whose candidate lists 1.2 widened.}
\label{fig:ka-changes}
\end{figure}

\begin{table}[t]
\centering
\footnotesize
\setlength{\tabcolsep}{3pt}
\begin{tabular}{@{}llr@{}}
\toprule
From & To & $n$ \\
\midrule
\multicolumn{3}{@{}l}{\emph{Run 1.1 to run 1.2}} \\
Logic & Mathematics & 87 \\
Business Law & Civil Law and Civil Procedure & 53 \\
Internal Medicine & General Surgery & 53 \\
Internal Medicine & Radiology and Medical Imaging & 51 \\
Physics & Radiology and Medical Imaging & 49 \\
Family and Community Med.\ and Public Health & Internal Medicine & 49 \\
Finance and Financial Systems & Mathematics & 47 \\
Nutrition & Gastronomy and Culinary Arts & 46 \\
Physics & Mathematics & 46 \\
Probability and Statistics & Mathematics & 44 \\
\midrule
\multicolumn{3}{@{}l}{\emph{Run 1.2 to run 1.4}} \\
General Surgery & Internal Medicine & 49 \\
Professional Ethics & Civil Law and Civil Procedure & 41 \\
Communication Studies & Portuguese Language & 37 \\
Mathematics & Logic & 34 \\
Internal Medicine & Family and Community Med.\ and Public Health & 30 \\
Literature & Portuguese Language & 26 \\
Programming, Algorithms and Data Structures & Logic & 25 \\
Constitutional Law & Administrative Law & 24 \\
Tax and Financial Law & Constitutional Law & 23 \\
Public Policy and Public Administration & Administrative Law & 22 \\
\bottomrule
\end{tabular}
\caption{The ten most frequent subject transitions between consecutive full runs on the 41,636 common questions.}
\label{tab:ka-transitions}
\end{table}

\paragraph{Consistency across runs.} The agreement between passes is
91.64\%, 92.14\% and 92.08\% in the three full runs and the adjudication
rate 8.36\%, 7.86\% and 7.92\%; the taxonomy versions did not change how
often the model is consistent with itself, only where it abstains and which
subject it chooses.

\paragraph{Reading.} A version that reduces abstentions and moves labels in
the exams whose lists it widened is doing what it was written to do; whether
the new labels are right is a separate question that these numbers do not
answer. The evidence files behind 1.2 and 1.3 are engineering reviews of the
model's output against the question text and were not produced by domain
experts; the 549 questions on new subjects and the largest transitions are the
natural sample for a human validation, which has not been done.

% -----------------------------------------------------------------------------
\subsection{Final label distribution}
\label{app:ka-distribution}

\paragraph{By level and macro area.} The two levels have different
profiles (Table~\ref{tab:ka-areas}). High-school questions are 38.05\%
Mathematics, Statistics and Logic (two subjects, driven by ENEM and OBI),
24.40\% Languages, 20.77\% Natural Sciences and 15.03\% Humanities, with 149
programming items from OBI and 22 Physical Education items as the only
labels outside those areas; three areas (Law, Economics, Social Sciences)
have no high-school question, since no vestibular list contains their
subjects. Undergraduate questions are 35.40\% Law, followed by Health
Sciences (12.94\%), Economics, Business and Accounting (10.70\%) and
Computing, Engineering, Architecture and Design (10.69\%). Over the whole
dataset Law holds 27.07\% of the questions and the smallest area, Social and
Applied Social Sciences, 5.94\% (2,472 questions), above the minimum of
2,000 per area that the evaluation protocol requires.

\begin{table}[t]
\centering
\footnotesize
\setlength{\tabcolsep}{4pt}
\begin{tabular}{@{}lrrr@{}}
\toprule
Macro area & Subjects & Questions & Share \\
\midrule
\multicolumn{4}{@{}l}{\emph{High school (9,789 questions)}} \\
Mathematics, Statistics, and Logic & 2 & 3,725 & 38.05 \\
Languages, Literature, Arts, and Communication & 6 & 2,389 & 24.40 \\
Natural, Earth, Agricultural and Veterinary Sciences & 3 & 2,033 & 20.77 \\
Humanities & 4 & 1,471 & 15.03 \\
Computing, Engineering, Architecture, and Design & 1 & 149 & 1.52 \\
Health Sciences & 1 & 22 & 0.22 \\
\midrule
\multicolumn{4}{@{}l}{\emph{Undergraduate (31,845 questions)}} \\
Law & 14 & 11,272 & 35.40 \\
Health Sciences & 15 & 4,122 & 12.94 \\
Economics, Business, and Accounting & 4 & 3,409 & 10.70 \\
Computing, Engineering, Architecture, and Design & 12 & 3,405 & 10.69 \\
Social and Applied Social Sciences & 9 & 2,472 & 7.76 \\
Humanities & 7 & 2,321 & 7.29 \\
Languages, Literature, Arts, and Communication & 7 & 1,975 & 6.20 \\
Natural, Earth, Agricultural and Veterinary Sciences & 7 & 1,534 & 4.82 \\
Mathematics, Statistics, and Logic & 3 & 1,335 & 4.19 \\
\bottomrule
\end{tabular}
\caption{Macro areas of the final dataset per academic level. \emph{Subjects}: distinct subjects used at that level; \emph{Share}: percentage of the level's questions. Over the whole dataset: Law 27.07\%; Mathematics, Statistics and Logic 12.15\%; Languages 10.48\%; Health Sciences 9.95\%; Humanities 9.11\%; Natural Sciences 8.57\%; Computing 8.54\%; Economics 8.19\%; Social Sciences 5.94\%.}
\label{tab:ka-areas}
\end{table}

\paragraph{By subject.} All 78 subjects are used (Figure~\ref{fig:ka-subjects},
Table~\ref{tab:ka-subjects}). Mathematics is the largest with 2,803
questions, followed by Civil Law and Civil Procedure (2,547), Logic (1,979),
Criminal Law and Criminal Procedure (1,839), Portuguese Language (1,755),
Constitutional Law (1,595), Labor Law and Labor Procedure (1,549) and
Business Administration (1,328). The tail is long: 14 subjects have fewer
than 100 questions and six fewer than 50, the smallest being Public Safety
and Physical Security (13), Human--Computer Interaction (25), Linguistics
(42), Social Security Law (42) and Human Rights (43). Only 17 subjects occur
at high-school level, and three of them (Spanish Language, Communication
Studies, Physical Education) with fewer than 30 questions there.

\begin{figure*}[p]
\centering
\includegraphics[width=\textwidth]{\KnowledgeAreaFigDir/subjects_by_level.pdf}
\caption{Questions per subject in the final dataset, stacked by academic level and grouped by macro area in taxonomy order. The scale is linear and shared by the two panels; the total is printed next to each bar.}
\label{fig:ka-subjects}
\end{figure*}

\begin{table*}[p]
\centering
\scriptsize
\setlength{\tabcolsep}{3pt}
\begin{tabular}{@{}lrrr@{}}
\toprule
Subject & HS & UG & Total \\
\midrule
\multicolumn{4}{@{}l}{\emph{Languages, Literature, Arts, and Communication}} \\
Portuguese Language & 1,054 & 701 & 1,755 \\
English Language & 735 & 319 & 1,054 \\
Spanish Language & 6 & 218 & 224 \\
Literature & 479 & 100 & 579 \\
Linguistics & -- & 42 & 42 \\
Arts & 87 & 121 & 208 \\
Communication Studies & 28 & 474 & 502 \\
\midrule
\multicolumn{4}{@{}l}{\emph{Mathematics, Statistics, and Logic}} \\
Mathematics & 1,874 & 929 & 2,803 \\
Logic & 1,851 & 128 & 1,979 \\
Probability and Statistics & -- & 278 & 278 \\
\midrule
\multicolumn{4}{@{}l}{\emph{Computing, Engineering, Architecture, and Design}} \\
Programming, Algorithms and Data Structures & 149 & 537 & 686 \\
Computer Systems and Architecture & -- & 254 & 254 \\
Theory of Computation & -- & 112 & 112 \\
Software Engineering and Databases & -- & 617 & 617 \\
Computer Networks and Distributed Systems & -- & 384 & 384 \\
Artificial Intelligence, Data Science and Computer Vision & -- & 110 & 110 \\
Information Security & -- & 137 & 137 \\
Engineering & -- & 747 & 747 \\
Architecture and Urban Planning & -- & 193 & 193 \\
Design & -- & 225 & 225 \\
Computer Graphics & -- & 64 & 64 \\
Human--Computer Interaction & -- & 25 & 25 \\
\midrule
\multicolumn{4}{@{}l}{\emph{Natural, Earth, Agricultural and Veterinary Sciences}} \\
Physics & 880 & 223 & 1,103 \\
Chemistry & 582 & 285 & 867 \\
Biology & 571 & 184 & 755 \\
Earth and Environmental Sciences & -- & 255 & 255 \\
Agronomy & -- & 307 & 307 \\
Veterinary Medicine & -- & 120 & 120 \\
Animal Science & -- & 160 & 160 \\
\midrule
\multicolumn{4}{@{}l}{\emph{Humanities}} \\
History & 645 & 237 & 882 \\
Geography & 519 & 222 & 741 \\
Philosophy & 114 & 330 & 444 \\
Sociology & 193 & 530 & 723 \\
Psychology & -- & 317 & 317 \\
Education and Pedagogy & -- & 632 & 632 \\
Theology and Religious Studies & -- & 53 & 53 \\
\bottomrule
\end{tabular}
\hfill
\begin{tabular}{@{}lrrr@{}}
\toprule
Subject & HS & UG & Total \\
\midrule
\multicolumn{4}{@{}l}{\emph{Social and Applied Social Sciences}} \\
Political Science & -- & 84 & 84 \\
International Relations & -- & 147 & 147 \\
Public Policy and Public Administration & -- & 370 & 370 \\
Social Work & -- & 127 & 127 \\
Information Science & -- & 509 & 509 \\
Tourism & -- & 125 & 125 \\
Professional Ethics & -- & 999 & 999 \\
Gastronomy and Culinary Arts & -- & 98 & 98 \\
Public Safety and Physical Security & -- & 13 & 13 \\
\midrule
\multicolumn{4}{@{}l}{\emph{Economics, Business, and Accounting}} \\
Economics & -- & 568 & 568 \\
Business Administration & -- & 1,328 & 1,328 \\
Finance and Financial Systems & -- & 404 & 404 \\
Accounting & -- & 1,109 & 1,109 \\
\midrule
\multicolumn{4}{@{}l}{\emph{Law}} \\
Constitutional Law & -- & 1,595 & 1,595 \\
Administrative Law & -- & 1,030 & 1,030 \\
Civil Law and Civil Procedure & -- & 2,547 & 2,547 \\
Criminal Law and Criminal Procedure & -- & 1,839 & 1,839 \\
Labor Law and Labor Procedure & -- & 1,549 & 1,549 \\
Business Law & -- & 935 & 935 \\
Tax and Financial Law & -- & 875 & 875 \\
International Law & -- & 289 & 289 \\
Environmental Law & -- & 223 & 223 \\
Consumer Law & -- & 147 & 147 \\
Electoral Law & -- & 44 & 44 \\
Social Security Law & -- & 42 & 42 \\
Human Rights & -- & 43 & 43 \\
Child and Adolescent Law & -- & 114 & 114 \\
\midrule
\multicolumn{4}{@{}l}{\emph{Health Sciences}} \\
Internal Medicine & -- & 1,126 & 1,126 \\
General Surgery & -- & 545 & 545 \\
Pediatrics & -- & 416 & 416 \\
Obstetrics and Gynecology & -- & 366 & 366 \\
Family and Community Medicine and Public Health & -- & 478 & 478 \\
Biomedicine & -- & 64 & 64 \\
Nursing & -- & 98 & 98 \\
Pharmacy & -- & 102 & 102 \\
Physical Therapy and Speech-Language Pathology & -- & 217 & 217 \\
Nutrition & -- & 163 & 163 \\
Dentistry & -- & 113 & 113 \\
Physical Education & 22 & 155 & 177 \\
Occupational Therapy & -- & 53 & 53 \\
Radiology and Medical Imaging & -- & 168 & 168 \\
Aesthetics and Cosmetology & -- & 58 & 58 \\
\bottomrule
\end{tabular}
\caption{The 78 subjects of taxonomy 1.4 and the questions assigned to each in the final dataset. \emph{HS}: high school; \emph{UG}: undergraduate; ``--'': no question of that level received the subject (the subject is not a candidate of any exam of that level).}
\label{tab:ka-subjects}
\end{table*}

\paragraph{By exam.} Table~\ref{tab:ka-exams} and Figure~\ref{fig:ka-exams}
show the labels of each exam. The exams with a declared scope stay inside it:
OAB is 91.3\% Law and 8.0\% Social and Applied Social Sciences (Professional
Ethics, a candidate of the bar exam), ENAM 97.1\% Law, REVALIDA 98.2\% and the
residency exams 94.3\% Health Sciences, CFCES 88.7\% Accounting, OBI 93.5\%
Mathematics, Statistics and Logic (80.65\% Logic) and POSCOMP 69.7\%
Computing and 30.3\% Mathematics, Statistics and Logic. The broad exams use
their lists widely: ENADE uses 77 of its 78 candidates and all nine areas
with Business Administration first at only 8.47\%; BNDES uses 35 of 36, CNU
39 of 51 and BACEN all 26. The vestibular exams spread over the school
subjects with no subject above 24\% in ENEM, FUVEST, BLUEX and COMVEST, while
the military academies split between languages and exact sciences (AFA:
Portuguese 26.45\%, Mathematics 25.71\%; IME: English 34.83\%, Mathematics
26.15\%). Since the exam and the edition are shown to the model, these
shares describe the labels, not an independent check of them.

\begin{table*}[t]
\centering
\footnotesize
\setlength{\tabcolsep}{4pt}
\begin{tabular}{@{}lrrrrlrlr@{}}
\toprule
Exam & $n$ & Cand. & Areas & Subjects & Most frequent subject & \% & Second subject & \% \\
\midrule
\multicolumn{9}{@{}l}{\emph{High school (9,789)}} \\
ENEM & 2,687 & 15 & 5 & 15 & Mathematics & 23.7 & Portuguese Language & 12.5 \\
OBI & 2,295 & 3 & 2 & 3 & Logic & 80.7 & Mathematics & 12.9 \\
FUVEST & 1,659 & 13 & 5 & 13 & History & 12.8 & Mathematics & 12.4 \\
AFA & 1,085 & 4 & 3 & 4 & Portuguese Language & 26.5 & Mathematics & 25.7 \\
IME & 956 & 5 & 3 & 5 & English Language & 34.8 & Mathematics & 26.2 \\
BLUEX & 656 & 11 & 4 & 11 & Mathematics & 18.0 & Literature & 17.4 \\
COMVEST & 451 & 11 & 4 & 11 & Mathematics & 19.7 & Geography & 14.0 \\
\midrule
\multicolumn{9}{@{}l}{\emph{Undergraduate (31,845)}} \\
ENADE & 10,785 & 78 & 9 & 77 & Business Administration & 8.5 & Engineering & 6.6 \\
OAB & 10,251 & 16 & 3 & 16 & Civil Law and Civil Procedure & 21.4 & Criminal Law and Criminal Procedure & 16.5 \\
BNDES & 4,048 & 36 & 8 & 35 & Information Science & 10.5 & Portuguese Language & 8.2 \\
Res.\ USP/Unicamp & 1,574 & 13 & 5 & 13 & Internal Medicine & 45.2 & General Surgery & 25.1 \\
POSCOMP & 1,302 & 12 & 2 & 12 & Programming, Algorithms and Data Structures & 23.0 & Mathematics & 22.3 \\
REVALIDA & 1,195 & 6 & 2 & 6 & Internal Medicine & 29.0 & Pediatrics & 21.6 \\
BACEN & 1,178 & 26 & 7 & 26 & Portuguese Language & 15.1 & Finance and Financial Systems & 7.4 \\
CFCES & 707 & 11 & 4 & 7 & Accounting & 88.7 & Business Law & 5.0 \\
CNU & 491 & 51 & 9 & 39 & Public Policy and Public Administration & 15.3 & Administrative Law & 11.0 \\
ENAM & 314 & 11 & 3 & 10 & Civil Law and Civil Procedure & 36.6 & Constitutional Law & 27.1 \\
\bottomrule
\end{tabular}
\caption{Labels per exam in the final dataset, sorted by size within each level. \emph{Cand.}: candidate subjects of the exam in taxonomy 1.4; \emph{Areas} and \emph{Subjects}: distinct macro areas and subjects assigned to the exam's questions; \emph{\%}: share of the exam's questions with the most frequent and the second most frequent subject.}
\label{tab:ka-exams}
\end{table*}

\begin{figure*}[t]
\centering
\includegraphics[width=\textwidth]{\KnowledgeAreaFigDir/macro_area_by_exam.pdf}
\caption{Share of each exam's questions assigned to every macro area (\%) in the final dataset. Exams are grouped by academic level and sorted by size; empty cells are macro areas without any question of that exam, and ``$<$1'' marks shares below one percent.}
\label{fig:ka-exams}
\end{figure*}

% -----------------------------------------------------------------------------
\subsection{Post-annotation}
\label{app:ka-post}

The eight abstentions of the final run were read by the authors and resolved
in one of two ways, recorded in two versioned lists that a separate command
applies to the run's export without touching the checkpoint. Six received a
manual subject, always one of the exam's candidates: in two cases the
adjudicator had rejected the subject of one pass claiming that it was not in
the list when it was (Spanish Language in a BNDES Spanish-comprehension item,
Professional Ethics in a bar-exam item on the statute of the profession), and
the pass's subject was restored; in three cases the question's discipline is
missing from the exam's list by design (no law subject in ENEM or COMVEST, no
administrative law in the residency exams) and the closest candidate was
chosen, with the alternative recorded; one is a boundary case between
General Surgery and Pediatrics in a head-and-neck surgery exam. These rows
carry the status \emph{manual}, a null confidence, a justification that
starts with the words ``manual assignment'' and the two model passes for
audit. Two questions were removed because their supporting chart or table is
missing from the extracted text; they join the 17 questions removed by the
dataset revision for the same family of defects (7 stems replaced by a
formatting instruction, 6 self-referential OBI items, 2 truncated stems and 2
items without their supporting text). The final dataset has 41,634 questions,
38,337 accepted, 3,291 adjudicated and 6 manual, no \texttt{UNCERTAIN}, and
was published as \texttt{bench-temp-2/mmlu-pt-knowledge-annotated} at
revision \texttt{fb9d509a}\ldots. The two adjudication errors contradict the
label--justification consistency rule of taxonomy 1.4 and are failures of the
model on single items rather than gaps of the taxonomy, which is why they were
corrected by hand instead of motivating a further version.

% -----------------------------------------------------------------------------
\subsection{Cost and reproducibility}
\label{app:ka-cost}

Table~\ref{tab:ka-cost} summarizes the final run. The 41,636 questions were
annotated on one NVIDIA B200 in 13,119~s (3\,h\,39\,min) from the start of
the job to the end of the export, of which 112~s were spent loading the
115.37~GiB of FP8 weights and initializing the engine, and 12,892~s were
spent inside generation batches (6,199~s for pass~1, 6,117~s for pass~2 and
576~s for the 3,299 adjudications). The run processed 311.2 million prompt
tokens and 6.2 million completion tokens: each of the two passes sends
148.4 million prompt tokens, 3,564 per question on average (2,020 at
high-school and 4,039 at undergraduate level, where the candidate lists are
longer; the longest prompt, an ENADE item, has 10,367 tokens), and receives about 72 completion tokens, the JSON object. Prefix
caching was disabled, so the prompt-token count is the full cost of
re-reading the rules and definitions for every question and an upper bound
for a re-run with caching. The earlier full runs cost 8,894~s and 162.4
million prompt tokens (1.1) and 11,301~s and 253.5 million (1.2): the
rules and definitions added in 1.2 raised the prompt tokens by 56\% and the
session time by 27\%, with the same number of batches. Reproducibility rests
on the run manifest, which records the source revision and its hash, the
model identifier and resolved revision, the taxonomy and alias hashes, the
prompt version and system message, the code hashes, the library versions and
the engine settings, and on the checkpoint, which stores the seed of every
request. Sampling at temperature~1 with per-request seeds reproduces the same
labels only on the same hardware, kernels and batch composition; on another
stack the labels should be expected to differ at the level of the within-run
disagreement.

\begin{table}[t]
\centering
\footnotesize
\setlength{\tabcolsep}{4pt}
\begin{tabular}{@{}lr@{}}
\toprule
Hardware and engine & \\
\midrule
GPU & 1$\times$ NVIDIA B200, 191.5~GB \\
Weights (FP8) & 115.37~GiB, loaded in 30.3~s \\
Engine & vLLM 0.25.0, XGrammar, CUDA graphs on \\
KV cache & 1,413,120 tokens \\
Context limit / output budget & 20,480 / 8,192 tokens \\
Batch size / concurrent sequences & 128 / 50 \\
Prefix caching & off \\
Job start to export & 13,119~s (3\,h\,39\,min) \\
Engine ready after & 112~s \\
\bottomrule
\end{tabular}

\medskip
\begin{tabular}{@{}lrrrr@{}}
\toprule
Stage & Requests & Prompt tokens & Compl.\ tokens & Batch time \\
\midrule
Pass 1 & 41,636 & 148,407,132 & 2,981,987 & 6,199~s \\
Pass 2 & 41,636 & 148,407,132 & 2,981,000 & 6,117~s \\
Adjudication & 3,299 & 14,429,317 & 255,387 & 576~s \\
Total & 86,571 & 311,243,581 & 6,218,374 & 12,892~s \\
\midrule
Run 1.1 (taxonomy 1.1) & 86,788 & 162,405,925 & 6,197,422 & 8,693~s \\
Run 1.2 (taxonomy 1.2) & 86,578 & 253,486,513 & 6,239,772 & 11,032~s \\
\bottomrule
\end{tabular}
\caption{Cost of the final run (top and middle) and totals of the two earlier full runs (bottom). \emph{Requests}: successful model requests; the six requests retried after a schema violation are excluded from the counts and included in the batch time. \emph{Batch time}: sum of the durations of the generation batches, deduplicated by batch; it excludes engine start-up and export.}
\label{tab:ka-cost}
\end{table}

% -----------------------------------------------------------------------------
\subsection{Limitations}
\label{app:ka-limitations}

\emph{No gold labels.} No question has a human-assigned subject, so nothing in
this appendix measures accuracy: agreement between passes, stability between
versions and the shares per exam measure the consistency of one model with
itself, and the audits that motivated the versions are engineering reviews of
the model's output, not expert annotation.
\emph{One model, not ablated.} Model, prompt template, decoding and
aggregation were fixed before the first run; the choice of model is motivated
by prior work, openness, size and hardware fit (\S\ref{app:ka-method}), not
by a comparison on this task, and six of the evaluated checkpoints belong to
the annotator's family, so a family-level affinity between the labels and
those models cannot be excluded.
\emph{Sampling noise.} The two passes disagree on about 8\% of the questions,
and the adjudicator is the same model; a change between versions on an exam is
interpretable only where it exceeds that level, and the 1.2-to-1.4 changes
are mostly within it.
\emph{Small targeted runs.} Versions 1.3 and 1.4 were tested on 23 and 15
questions before the third full run; their effect on the rest of the corpus
is only visible through the full run, where it is confounded with sampling.
\emph{Metadata in the prompt.} The exam and the edition are shown to the
model as context, so the per-exam shares of \S\ref{app:ka-distribution} are
not an independent check of the labels.
\emph{Uncalibrated confidence.} The reported confidence is a number the model
writes; 21 of the 23 abstentions of run 1.2 carried a confidence of 0.90 or
more.
\emph{One label per question.} Interdisciplinary questions receive one
subject, with an alternative stored but unused, and the candidate lists of the
vestibular and residency exams do not contain every discipline their
questions touch, as three of the six manual assignments show.
\emph{Constrained decoding.} The grammar guarantees an admissible label, not a
correct one, and format restrictions can affect the quality of a model's
reasoning \citep{tam2024letme}; the run constrains only the final object,
with the reasoning mode off.
\emph{Project-specific taxonomy.} The subjects were written for these 17
exams and do not follow an official classification of fields; per-subject
results are not directly comparable to other benchmarks, and 14 subjects have
fewer than 100 questions.
\emph{Changing source.} Run 1.4 annotated a revised corpus, so 17 questions of
the earlier runs are absent from the comparison, and the final dataset drops
two more.
